\nonstopmode{}
\documentclass[a4paper]{book}
\usepackage[times,inconsolata,hyper]{Rd}
\usepackage{makeidx}
\makeatletter\@ifl@t@r\fmtversion{2018/04/01}{}{\usepackage[utf8]{inputenc}}\makeatother
% \usepackage{graphicx} % @USE GRAPHICX@
\makeindex{}
\begin{document}
\chapter*{}
\begin{center}
{\textbf{\huge Package `nutricalc'}}
\par\bigskip{\large \today}
\end{center}
\ifthenelse{\boolean{Rd@use@hyper}}{\hypersetup{pdftitle = {nutricalc: Nutrient Solution Formulation in Horticultural Sciences}}}{}
\ifthenelse{\boolean{Rd@use@hyper}}{\hypersetup{pdfauthor = {Dilan Olmo Sakar}}}{}
\begin{description}
\raggedright{}
\item[Type]\AsIs{Package}
\item[Title]\AsIs{Nutrient Solution Formulation in Horticultural Sciences}
\item[Version]\AsIs{0.1.0}
\item[Description]\AsIs{Provides tools for parsing chemical formulas, calculating molar masses,
and optimizing fertilizer mixtures to meet specified nutrient targets using 
non-negative least squares. Designed to help agronomists and researchers formulate 
custom nutrient solutions efficiently.}
\item[License]\AsIs{GPL-3}
\item[Encoding]\AsIs{UTF-8}
\item[LazyData]\AsIs{true}
\item[Depends]\AsIs{R (>= 3.5)}
\item[Imports]\AsIs{stringr,
nnls,
shiny,
bslib,
Ternary,
xml2}
\item[RoxygenNote]\AsIs{7.3.2}
\end{description}
\Rdcontents{Contents}
\HeaderA{assign\_salts\_bak}{Assign salts to A, B, and Micro tanks based on composition and solver result}{assign.Rul.salts.Rul.bak}
%
\begin{Description}
This helper takes the result from \code{optimize\_nutrients()} and the
corresponding \code{nutrient\_matrix} (compounds x nutrients) and returns a
BAK-style grouping:
\end{Description}
%
\begin{Usage}
\begin{verbatim}
assign_salts_bak(
  res,
  nutrient_matrix,
  micro_nutrients = c("Fe", "Mn", "Zn", "B", "Cu", "Mo", "Si"),
  calcium_nutrient = "Ca",
  phosphorus_nutrient = "P",
  sulfur_nutrient = "S",
  balance_nutrients = c("NO3_N", "NH4_N", "P", "K", "Ca", "Mg", "S", "Na"),
  drop_zero = TRUE,
  tol = 0
)
\end{verbatim}
\end{Usage}
%
\begin{Arguments}
\begin{ldescription}
\item[\code{res}] A \code{nutrient\_optimization\_result} as returned by
\code{optimize\_nutrients()}.

\item[\code{nutrient\_matrix}] Matrix [compound x nutrient] used in the optimization.
Row names must match the salt names (formulas).

\item[\code{micro\_nutrients}] Character vector of column names that count as
micronutrients. Defaults to \code{c("Fe", "Mn", "Zn", "B", "Cu", "Mo")}.

\item[\code{calcium\_nutrient}] Column name for calcium (default: \code{"Ca"}).

\item[\code{phosphorus\_nutrient}] Column name for phosphorus (default: \code{"P"}).

\item[\code{sulfur\_nutrient}] Column name for sulfur (default: \code{"S"}).

\item[\code{balance\_nutrients}] Character vector of nutrient columns used to define
the "load" when splitting AB salts between A and B. Defaults to
\code{c("NO3\_N","NH4\_N","P","K","Ca","Mg","S","Na")}.

\item[\code{drop\_zero}] Logical, if \code{TRUE} (default), remove salts with
effectively zero mmol/L from the returned \code{table} and \code{by\_tank}.
The \code{assignments} vector always includes all salts.

\item[\code{tol}] Numeric tolerance for deciding what counts as "zero" mmol/L.
\end{ldescription}
\end{Arguments}
%
\begin{Details}
- First, salts are classified purely by composition:
* Micro  = any micro nutrient present (e.g. Fe, Mn, Zn, B, Cu, Mo)
* A      = contains Ca (and no micro)
* B      = contains P or S (and no micro / Ca)
* AB     = none of the above (can go to A or B)

- Then the AB salts are assigned to A or B so that A and B have a similar
overall "load" (sum of delivered mmol of selected nutrients).
\end{Details}
%
\begin{Value}
A list with:
\begin{ldescription}
\item[\code{assignments}] Named character vector, salt -> tank ("A","B","Micro").
\item[\code{table}] Data frame with columns \code{salt}, \code{tank}, \code{mmol\_l}
(filtered to non-zero if \code{drop\_zero = TRUE}).
\item[\code{by\_tank}] Named list of data frames (A, B, Micro, possibly others),
each subset of \code{table}.
\end{ldescription}
\end{Value}
\HeaderA{CANONICAL\_NUTRIENTS}{Canonical nutrient keys for targets and UI}{CANONICAL.Rul.NUTRIENTS}
\keyword{datasets}{CANONICAL\_NUTRIENTS}
%
\begin{Description}
Defines the stable nutrient order used across recipes, UI inputs, and target
vectors. This is intentionally independent from the full `nutrient\_matrix`,
which may include additional columns for chelate ligands or other species.
\end{Description}
%
\begin{Usage}
\begin{verbatim}
CANONICAL_NUTRIENTS
\end{verbatim}
\end{Usage}
%
\begin{Format}
A character vector of nutrient keys.
\end{Format}
\HeaderA{compute\_molar\_mass}{Compute molar mass of a chemical formula}{compute.Rul.molar.Rul.mass}
%
\begin{Description}
Supports formulas like "NaCl", "Ca(NO3)2", "CuSO4·5H2O", "MgN2O6·6H2O"
\end{Description}
%
\begin{Usage}
\begin{verbatim}
compute_molar_mass(formula)
\end{verbatim}
\end{Usage}
%
\begin{Arguments}
\begin{ldescription}
\item[\code{formula}] Chemical formula string
\end{ldescription}
\end{Arguments}
%
\begin{Value}
Molar mass in g/mol
\end{Value}
\HeaderA{compute\_percentage\_matrix}{Compute a matrix of mass percentages of nutrients in each fertilizer}{compute.Rul.percentage.Rul.matrix}
%
\begin{Description}
Compute a matrix of mass percentages of nutrients in each fertilizer
\end{Description}
%
\begin{Usage}
\begin{verbatim}
compute_percentage_matrix(nutrient_matrix)
\end{verbatim}
\end{Usage}
%
\begin{Arguments}
\begin{ldescription}
\item[\code{nutrient\_matrix}] A [fertilizer x nutrient] matrix in mmol/mol
\end{ldescription}
\end{Arguments}
%
\begin{Value}
A matrix of nutrient mass percentages per fertilizer
\end{Value}
\HeaderA{converte}{Convert amount of one compound to another using molar mass and element counts}{converte}
%
\begin{Description}
This function computes the conversion factor between two chemical formulas or elements
based on their molar masses and the count of shared atoms. It adjusts for stoichiometry
where applicable (e.g., converting from "SO4" to "S").
\end{Description}
%
\begin{Usage}
\begin{verbatim}
converte(from, to, value = 1, element = NULL)
\end{verbatim}
\end{Usage}
%
\begin{Arguments}
\begin{ldescription}
\item[\code{from}] A character string representing the source formula or element (e.g., "Ca", "NaCl", "SO4").

\item[\code{to}] A character string representing the target formula or element.

\item[\code{value}] A numeric value to convert (default is 1). Represents the quantity of `from`.

\item[\code{element}] Optional element symbol to anchor the stoichiometric conversion when
multiple elements are shared between `from` and `to`. If omitted and more than one
element is shared, an error is raised to avoid ambiguous conversions.
\end{ldescription}
\end{Arguments}
%
\begin{Value}
A numeric value representing the equivalent amount of `to`, adjusted by molar mass and atom ratio.
\end{Value}
%
\begin{Examples}
\begin{ExampleCode}
converte("SO4", "S")        # Convert sulfate to elemental sulfur
converte("Na", "NaCl")      # Convert sodium to sodium chloride
converte("P2O5", "P")       # Convert phosphorus pentoxide to elemental phosphorus
converte("Ca", "CaCl2\u00b72H2O", 10)  # Convert 10 mmol of Ca to CaCl2\u00b72H2O
converte("(NH4)2HPO4", "NH4H2PO4", element = "P")  # Disambiguate using P as the anchor

\end{ExampleCode}
\end{Examples}
\HeaderA{convert\_units}{Convert nutrient concentrations between mmol/L and mg/L}{convert.Rul.units}
%
\begin{Description}
Converts numeric nutrient values (single element or named vector)
between mmol/L and mg/L using molar masses in `element\_molar\_mass\_df`.
\end{Description}
%
\begin{Usage}
\begin{verbatim}
convert_units(x, element = NULL, to = c("mg/L", "mmol/L"))
\end{verbatim}
\end{Usage}
%
\begin{Arguments}
\begin{ldescription}
\item[\code{x}] Numeric scalar or named numeric vector of nutrient values.
If named, names should match those in `nutrient\_element\_map`.

\item[\code{element}] Optional element symbol (e.g. "N", "P") when `x` is a scalar.

\item[\code{to}] Target unit: "mg/L" or "mmol/L".
\end{ldescription}
\end{Arguments}
%
\begin{Value}
Numeric vector (same length/names as `x`) in the target unit.
\end{Value}
%
\begin{Examples}
\begin{ExampleCode}
convert_units(c(NO3_N = 15, P = 1.5), to = "mg/L")
convert_units(c(NO3_N = 210, P = 46), to = "mmol/L")
\end{ExampleCode}
\end{Examples}
\HeaderA{drop\_tiny}{Zero-out tiny floating point values}{drop.Rul.tiny}
%
\begin{Description}
Utility helper to clamp extremely small values to zero, avoiding noisy output.
\end{Description}
%
\begin{Usage}
\begin{verbatim}
drop_tiny(x, tol = 1e-09)
\end{verbatim}
\end{Usage}
%
\begin{Arguments}
\begin{ldescription}
\item[\code{x}] Numeric vector or matrix.

\item[\code{tol}] Values with absolute magnitude below this tolerance are set to zero.
\end{ldescription}
\end{Arguments}
%
\begin{Value}
The input with near-zero entries replaced by zeros.
\end{Value}
%
\begin{Examples}
\begin{ExampleCode}
drop_tiny(c(1, 1e-10, -1e-12))
drop_tiny(matrix(c(0.5, 1e-12), nrow = 1))
\end{ExampleCode}
\end{Examples}
\HeaderA{ec\_from\_ph}{Compute EC from achieved composition and pH speciation}{ec.Rul.from.Rul.ph}
%
\begin{Description}
Estimates conductivity using an ion-transport summation based on diffusion
coefficients and activity coefficients.
\end{Description}
%
\begin{Usage}
\begin{verbatim}
ec_from_ph(
  achieved,
  ph_res,
  fe_state = c("Fe2+", "Fe3+"),
  t_C = 25,
  gamma = NULL,
  gamma_model = 1,
  A_DH_25 = 0.5085
)
\end{verbatim}
\end{Usage}
%
\begin{Arguments}
\begin{ldescription}
\item[\code{achieved}] Named numeric vector (mmol/L) or data.frame with columns
Nutrient and Achieved.

\item[\code{ph\_res}] Result list from `ph\_from\_achieved()` containing `species\_mM`.

\item[\code{fe\_state}] One of "Fe2+" or "Fe3+". When "Fe3+", all Fe2+ is treated
as Fe3+ for EC purposes.

\item[\code{t\_C}] Temperature in degrees C.

\item[\code{gamma}] Optional named numeric vector of activity coefficients (per ion).

\item[\code{gamma\_model}] Activity model selector: 1/"davies\_25C" (Davies) or 2/"debye\_huckel\_25C" (Debye-Huckel). Default 1.

\item[\code{A\_DH\_25}] Debye-Huckel A constant at 25 C (also used as A in Davies).
\end{ldescription}
\end{Arguments}
%
\begin{Details}
Units:
- c is in mmol/L (mM); numerically identical to mol/m\textasciicircum{}3, used directly as such.
- D is in 10\textasciicircum{}-5 cm\textasciicircum{}2/s and converted to m\textasciicircum{}2/s by multiplying 1e-9.
- Output: EC in mS/cm and uS/cm.

gamma\_model:
1 = Davies (default)
2 = Debye-Hückel

Override precedence:
model baseline -> ph\_res\$gamma\_ions -> gamma argument (highest priority)
\end{Details}
%
\begin{Value}
A list with EC\_mS\_cm, EC\_uS\_cm, and per-ion contributions.
\end{Value}
\HeaderA{element\_molar\_mass\_df}{Atomic molar masses used for formula parsing}{element.Rul.molar.Rul.mass.Rul.df}
\keyword{datasets}{element\_molar\_mass\_df}
%
\begin{Description}
A lookup table of elemental (atomic) molar masses in grams per mole (g/mol).
Used by molar-mass and unit-conversion helpers (e.g., when computing the molar
mass of a chemical formula from its elemental composition).
\end{Description}
%
\begin{Usage}
\begin{verbatim}
data(element_molar_mass_df)
\end{verbatim}
\end{Usage}
%
\begin{Format}
A data frame with 23 rows and 2 columns:
\begin{description}

\item[Element] Element symbol (character), e.g. \code{"C"}, \code{"O"}, \code{"Fe"}.
\item[MolarMass\_g\_per\_mol] Atomic molar mass in g/mol (numeric).

\end{description}

\end{Format}
\HeaderA{fetch\_bwb\_mittelwert}{Fetch BWB "In aller Tiefe" Mittelwerte by postal code}{fetch.Rul.bwb.Rul.mittelwert}
%
\begin{Description}
Retrieves Berliner Wasserbetriebe (BWB) analysis data for a given Berlin
postal code and returns the Mittelwert values (first available entry per
parameter) for selected parameters as mmol/L.
\end{Description}
%
\begin{Usage}
\begin{verbatim}
fetch_bwb_mittelwert(plz)
\end{verbatim}
\end{Usage}
%
\begin{Arguments}
\begin{ldescription}
\item[\code{plz}] Character or numeric vector of length one with a five digit
German postal code.
\end{ldescription}
\end{Arguments}
%
\begin{Value}
A named numeric vector of length 18 in mmol/L with names
`c("NO3\_N","NH4\_N","P","K","Ca","Mg","S","Na","Cl",
"Fe","Mn","Zn","B","Cu","Mo","Si","KS4\_3","KB8\_2")`, limited to the decimal precision
provided by the source values. `KS4\_3` is the acid capacity (alkalinity, KS 4.3);
`KB8\_2` is the base capacity (acidity, KS 8.2) and is treated as dissolved CO2(aq).
\end{Value}
\HeaderA{ion\_diffusion\_25C}{Lookup table of ion diffusion coefficients at 25 C}{ion.Rul.diffusion.Rul.25C}
%
\begin{Description}
Returns a data frame with ions, charge, and diffusion coefficients (D) at 25 C.
Values are expressed as 10\textasciicircum{}-5 cm\textasciicircum{}2/s.
\end{Description}
%
\begin{Usage}
\begin{verbatim}
ion_diffusion_25C()
\end{verbatim}
\end{Usage}
%
\begin{Value}
A data.frame with columns: ion, z, D\_1e5\_cm2\_s.
\end{Value}
\HeaderA{nutrient\_matrix}{Canonical fertilizer-to-nutrient stoichiometry matrix}{nutrient.Rul.matrix}
\keyword{datasets}{nutrient\_matrix}
%
\begin{Description}
A compound-by-nutrient matrix used by the nutrient solvers. Rows are fertilizer
compounds (salts/acids/bases) and columns are nutrients. Each entry is the
amount of nutrient delivered by 1 mmol of the compound (mmol nutrient per mmol
compound).
\end{Description}
%
\begin{Usage}
\begin{verbatim}
data(nutrient_matrix)
\end{verbatim}
\end{Usage}
%
\begin{Format}
A numeric matrix with n rows (compounds) and 20 columns (nutrients).
Row names are chemical formulas; column names are nutrient keys:
`NO3\_N`, `NH4\_N`, `P`, `K`, `Ca`, `Mg`, `S`, `Na`, `Cl`, `Fe`, `Mn`, `Zn`,
`B`, `Cu`, `Mo`, `Si`, `EDTA`, `DTPA`, `EDDHA`, `HBED`.
\end{Format}
%
\begin{Details}
The matrix entries are stoichiometric coefficients in mmol/mmol. For example,
calcium nitrate (`Ca(NO3)2`) contributes 2 mmol `NO3\_N` and 1 mmol `Ca` per
1 mmol of salt.
\end{Details}
\HeaderA{optimize\_nutrients}{Optimize Nutrient Solution Using NNLS}{optimize.Rul.nutrients}
%
\begin{Description}
Solves for optimal fertilizer amounts to meet target nutrient concentrations.
Expects a compound-by-nutrient matrix: rows = salts/compounds, columns = nutrients.
\end{Description}
%
\begin{Usage}
\begin{verbatim}
optimize_nutrients(nutrient_matrix, target, importance)
\end{verbatim}
\end{Usage}
%
\begin{Arguments}
\begin{ldescription}
\item[\code{nutrient\_matrix}] Matrix [compound x nutrient] with entries in mol of nutrient per mmol of compound
(i.e., how many mmol of nutrient you get from 1 mmol of the compound).

\item[\code{target}] Named numeric vector of target concentrations (mmol/L).

\item[\code{importance}] Named numeric vector (-2 to +2) specifying nutrient importance (weights).
\end{ldescription}
\end{Arguments}
%
\begin{Value}
An object of class 'nutrient\_optimization\_result' with elements:
\begin{itemize}

\item{} amounts - named numeric (mmol/L) of each compound (same order as rownames of input)
\item{} achieved - named numeric (mmol/L) delivered per nutrient
\item{} target - named numeric (mmol/L) target per nutrient
\item{} abs\_error - achieved - target (mmol/L)
\item{} percent\_error - 100 * abs\_error / target (NA when target == 0)
\item{} squared\_error - sum(abs\_error\textasciicircum{}2)
\item{} rel\_squared\_error - sum((abs\_error/target)\textasciicircum{}2, na.rm = TRUE)

\end{itemize}

\end{Value}
\HeaderA{percentage\_matrix}{Nutrient mass percentage matrix}{percentage.Rul.matrix}
%
\begin{Description}
Mass percentages of nutrients in each fertilizer derived from `nutrient\_matrix`.
\end{Description}
%
\begin{Usage}
\begin{verbatim}
data(percentage_matrix)
\end{verbatim}
\end{Usage}
%
\begin{Format}
A numeric matrix with rows as fertilizer formulas and columns as nutrients
(NO3\_N, NH4\_N, P, K, Ca, Mg, S, Na, Cl, Fe, Mn, Zn, B, Cu, Mo, Si). Each value
is the mass percentage of the nutrient in the fertilizer.
\end{Format}
\HeaderA{ph\_from\_achieved}{Compute pH of a nutrient solution from achieved composition}{ph.Rul.from.Rul.achieved}
%
\begin{Description}
Uses charge balance with acid/base speciation and activity corrections
to estimate pH at 25 C. Assumes no carbonate alkalinity, no complexation, and
no precipitation.
\end{Description}
%
\begin{Usage}
\begin{verbatim}
ph_from_achieved(
  achieved,
  temp_C = 25,
  phc_bracket = c(3, 9),
  tol = 1e-09,
  max_iter = 200,
  inner_max_iter = 50,
  debug = FALSE,
  chelate_z = c(EDTA = -4, DTPA = -5, EDDHA = -4, HBED = -4),
  gamma_model = 1
)
\end{verbatim}
\end{Usage}
%
\begin{Arguments}
\begin{ldescription}
\item[\code{achieved}] Named numeric vector of mmol/L (e.g., solver result `x\$achieved`).
If present, `CO2\_aq` is treated as dissolved free CO2(aq) (mmol/L), mapped
from BWB "Basenkapazitaet KB 8,2" / KS 8.2 and used to fix carbonate speciation.

\item[\code{temp\_C}] Temperature in Celsius (only 25 C supported in this version).

\item[\code{phc\_bracket}] Numeric length-2 bracket for pHc = -log10([H+]) search.

\item[\code{tol}] Root tolerance in meq/L.

\item[\code{max\_iter}] Maximum iterations for the outer bisection.

\item[\code{inner\_max\_iter}] Maximum iterations for the ionic strength fixed point.

\item[\code{debug}] Logical; when TRUE prints basic bracket diagnostics.

\item[\code{chelate\_z}] Named numeric vector of chelate charges to include as fixed
anions. Values are net charges (negative) for ligand totals in mmol/L.

\item[\code{gamma\_model}] Integer 1 (Davies) or 2 (Debye–Hückel). Default 1.
\end{ldescription}
\end{Arguments}
%
\begin{Details}
gamma\_model:
1 = Davies (default)
2 = Debye–Hückel limiting law (DH)
\end{Details}
%
\begin{Value}
A list with pH, ionic strength, activity coefficients, species
distribution, and charge balance diagnostics.
\end{Value}
\HeaderA{plot\_ternary}{Steiner-style ternary plot (Target vs Achieved) with constant-component drops}{plot.Rul.ternary}
%
\begin{Description}
Steiner-style ternary plot (Target vs Achieved) with constant-component drops
\end{Description}
%
\begin{Usage}
\begin{verbatim}
plot_ternary(
  result,
  keys,
  title = "Ternary (ratios)",
  show_legend = TRUE,
  step_percent = 10,
  lab_cex = 0.95,
  axis_cex = 0.85,
  padding = 0.08,
  col_target = "black",
  col_achieved = "firebrick",
  pch_target = 16,
  pch_achieved = 17,
  cex_points = 1.25,
  show_link = FALSE,
  lty_link = 2,
  lwd_link = 1,
  legend_pos = "topright",
  show_axis_drops = TRUE,
  drop_col = "grey35",
  drop_lwd = 1,
  drop_lty = 1,
  show_ratio = TRUE,
  ratio_format = "percent",
  ratio_digits = 1
)
\end{verbatim}
\end{Usage}
%
\begin{Arguments}
\begin{ldescription}
\item[\code{result}] List with numeric named vectors `target` and `achieved`.

\item[\code{keys}] Character(3) giving axis order, e.g. c("K","Ca","Mg").

\item[\code{title}] Plot title.

\item[\code{show\_legend}] Logical, show legend (default TRUE).

\item[\code{step\_percent}] Integer in [1, 50]; major grid step as percent. Default 10.

\item[\code{lab\_cex}, \code{axis\_cex}, \code{padding}] Plot layout controls.

\item[\code{col\_target}, \code{col\_achieved}] Colors.

\item[\code{pch\_target}, \code{pch\_achieved}] Point symbols.

\item[\code{cex\_points}] Point size.

\item[\code{show\_link}] Logical, draw connector between Target and Achieved.

\item[\code{lty\_link}, \code{lwd\_link}] Style for connector if shown.

\item[\code{legend\_pos}] Legend position.

\item[\code{show\_axis\_drops}] Logical, draw constant-component drops from Achieved to sides.

\item[\code{drop\_col}, \code{drop\_lwd}, \code{drop\_lty}] Style for drop lines.

\item[\code{show\_ratio}] Logical, include "(A, B; C)" values in legend (default TRUE).

\item[\code{ratio\_format}] "percent" (default) or "proportion".

\item[\code{ratio\_digits}] Digits for numeric formatting.
\end{ldescription}
\end{Arguments}
%
\begin{Value}
Invisibly returns NULL (draws the plot).
\end{Value}
\HeaderA{print.bwb\_mittelwert}{Pretty-print BWB Mittelwerte}{print.bwb.Rul.mittelwert}
%
\begin{Description}
Pretty-print BWB Mittelwerte
\end{Description}
%
\begin{Usage}
\begin{verbatim}
## S3 method for class 'bwb_mittelwert'
print(x, ...)
\end{verbatim}
\end{Usage}
%
\begin{Arguments}
\begin{ldescription}
\item[\code{x}] A `bwb\_mittelwert` numeric vector.

\item[\code{...}] Additional arguments passed to `print`.
\end{ldescription}
\end{Arguments}
%
\begin{Value}
The input `bwb\_mittelwert` object, invisibly.
\end{Value}
\HeaderA{print.nutrient\_optimization\_result}{Pretty-print the optimization result}{print.nutrient.Rul.optimization.Rul.result}
%
\begin{Description}
Shows per-liter amounts; if `vol > 1`, also shows a dynamic column "g for <vol> L".
Uses `compute\_molar\_mass(formula)` (g/mol) when available; if not found, prints without mass columns.
\end{Description}
%
\begin{Usage}
\begin{verbatim}
## S3 method for class 'nutrient_optimization_result'
print(x, vol = 1, ...)
\end{verbatim}
\end{Usage}
%
\begin{Arguments}
\begin{ldescription}
\item[\code{x}] A 'nutrient\_optimization\_result' object.

\item[\code{vol}] Total solution volume in liters (must be > 0). If > 1, a "g for <vol> L" column is added.

\item[\code{...}] Unused.
\end{ldescription}
\end{Arguments}
\HeaderA{recipes}{Nutrient solution recipes used by NutriCalc}{recipes}
\keyword{datasets}{recipes}
%
\begin{Description}
A named list of recipe definitions. Each element is a list with:
\begin{description}

\item[name] Display name
\item[category] Category such as "Olericulture", "Fruticulture", ...
\item[notes] Free text notes / citation
\item[unit] Either "mmol/L" or "mg/L"
\item[salts] TRUE (all salts) or a character vector of salt formulas
\item[targets] Named numeric vector of nutrient targets

\end{description}

\end{Description}
%
\begin{Usage}
\begin{verbatim}
recipes
\end{verbatim}
\end{Usage}
%
\begin{Format}
A named list of recipe lists.
\end{Format}
\HeaderA{recipe\_summary}{Create a summary table of nutrient concentrations in both units}{recipe.Rul.summary}
%
\begin{Description}
Given a named vector of nutrient concentrations and its unit,
returns a data frame showing each nutrient's value in mmol/L and mg/L.
\end{Description}
%
\begin{Usage}
\begin{verbatim}
recipe_summary(recipe, unit = c("mmol/L", "mg/L"))
\end{verbatim}
\end{Usage}
%
\begin{Arguments}
\begin{ldescription}
\item[\code{recipe}] Named numeric vector of nutrient concentrations (e.g., NO3\_N, P, K, ...).

\item[\code{unit}] Current unit of the values: "mmol/L" or "mg/L".
\end{ldescription}
\end{Arguments}
%
\begin{Value}
A data frame with columns `Nutrient`, `Element`, `mmol/L`, `mg/L`.
\end{Value}
%
\begin{Examples}
\begin{ExampleCode}
recipe_summary(c(NO3_N = 15, P = 1.5), unit = "mmol/L")
recipe_summary(c(NO3_N = 210, P = 46), unit = "mg/L")
\end{ExampleCode}
\end{Examples}
\HeaderA{run\_app}{Launch the NutriCalc Shiny App}{run.Rul.app}
%
\begin{Description}
Launch the NutriCalc Shiny App
\end{Description}
%
\begin{Usage}
\begin{verbatim}
run_app(...)
\end{verbatim}
\end{Usage}
%
\begin{Arguments}
\begin{ldescription}
\item[\code{...}] Arguments passed to [shiny::runApp()].
\end{ldescription}
\end{Arguments}
\HeaderA{salt\_info}{Metadata for fertilizer compounds}{salt.Rul.info}
\keyword{datasets}{salt\_info}
%
\begin{Description}
Descriptive metadata for each compound in \code{\LinkA{nutrient\_matrix}{nutrient.Rul.matrix}},
including a human-readable category name and whether the compound is treated
as an acid/base in the two-stage solver.
\end{Description}
%
\begin{Usage}
\begin{verbatim}
data(salt_info)
\end{verbatim}
\end{Usage}
%
\begin{Format}
A data frame with n rows and 3 columns:
\begin{description}

\item[salt] Chemical formula key matching \code{rownames(nutrient\_matrix)}.
\item[category] Human-readable compound name/category.
\item[is\_acid\_base] Logical flag; \code{TRUE} for acids/bases.

\end{description}

\end{Format}
\HeaderA{strip\_negative\_zero}{Remove negative zeros introduced by rounding}{strip.Rul.negative.Rul.zero}
%
\begin{Description}
Ensures rounded values like -0.00 are displayed as 0.
\end{Description}
%
\begin{Usage}
\begin{verbatim}
strip_negative_zero(x)
\end{verbatim}
\end{Usage}
%
\begin{Arguments}
\begin{ldescription}
\item[\code{x}] Numeric vector or matrix.
\end{ldescription}
\end{Arguments}
\HeaderA{two\_stage\_optimize\_nutrients}{Two-stage nutrient optimization separating neutral salts from acids/bases}{two.Rul.stage.Rul.optimize.Rul.nutrients}
%
\begin{Description}
Runs an initial NNLS optimization on neutral salts, then uses acids/bases to
fill remaining deficiencies (when they can affect the off-target nutrients).
\end{Description}
%
\begin{Usage}
\begin{verbatim}
two_stage_optimize_nutrients(
  nutrient_matrix,
  target,
  importance,
  tol_abs = 0.01,
  tol_rel = 0.01
)
\end{verbatim}
\end{Usage}
%
\begin{Arguments}
\begin{ldescription}
\item[\code{nutrient\_matrix}] Matrix [compound x nutrient] with entries in mol of nutrient per mmol of compound.

\item[\code{target}] Named numeric vector of target concentrations (mmol/L).

\item[\code{importance}] Named numeric vector (-2 to +2) specifying nutrient importance (weights).

\item[\code{tol\_abs}] Absolute tolerance (mmol/L) below which residuals are ignored for acid/base correction.

\item[\code{tol\_rel}] Relative tolerance (fraction of target) below which residuals are ignored for acid/base correction.
\end{ldescription}
\end{Arguments}
%
\begin{Value}
A 'nutrient\_optimization\_result' object.
\end{Value}
\printindex{}
\end{document}
